\section{Conclusion}
\label{sec:conclusion}

We set out to study whether a task-grounded uncertainty framework can improve
LLM-based reranking under a controlled same-candidate protocol, and our main
contribution is the protocol itself: every method ranks the same 101 candidates
per event for 10{,}000 users, with a shared Qwen3-8B backbone, a shared
importer, identical metrics, exact coverage checks, and paired Holm-corrected
bootstrap tests. Under this protocol, our central finding is that a
task-grounded pointwise LLM relevance posterior---the model's own per-candidate
relevance probability used directly as the ranking score---ranks first against
eight official-code-level baselines in six of the eight Amazon domains we tested
(Books, Electronics, Sports, Toys, Home, Tools), improving NDCG@10 by $+21.6\%$
to $+53.2\%$ over the strongest baseline and leading on essentially all seven
HR/NDCG/MRR metrics in those six. In the other two domains it is competitive but
not first---rank~2 on Beauty ($-11\%$ NDCG@10) and rank~5 on Movies
($-24\%$)---which we report rather than omit. On the four domains with full
per-event signal rows all 56 per-domain paired tests are positive and
Holm-significant. This is a scoped same-candidate reranking claim, not a
full-catalog recommender SOTA result.

The diagnostic modules locate this signal precisely and yield a rigorous
negative result. Leave-one-component-out ablations and a raw-probability
attribution show that the uncertainty decomposition and the risk-adjusted
ranking family do not improve same-candidate ranking over the posterior they are
built on: the test-best risk exponent is zero in all four domains, a
confidence-only variant matches or exceeds the full configuration, and the
individual uncertainty terms are inert or mildly harmful. The working ranking
signal lies in the pointwise posterior itself, not in the uncertainty machinery.
The event-level uncertainty signal does stratify ranking reliability, but we
report this only as a caveated, partly mechanical descriptor, not as evidence
that uncertainty improves ranking. We present the calibrated-uncertainty and
risk-adjusted family as an analyzed and characterized design space, and leave
its use as a reliability descriptor or for downstream abstention to future work.
